\textbf{Program 2: Matrix operations in one program.}

\textbf{File:} \texttt{matrix\_operations.cpp}.

\textbf{Chapter connection:} Arrays and Two-Dimensional Arrays: row/column indexing.

Use the same 3 by 3 input for row sums, column sums, transpose, maximum, diagonal sum, and symmetry. Each function states its traversal through its loop bounds. main also supplies symmetric and nonsymmetric examples.

\begin{adslabfigure}[How this program runs]
\begin{adslabgraphic}
\begin{tikzpicture}[font=\small,execute at begin node={\hyphenpenalty=10000\exhyphenpenalty=10000}]
\node[utpstep,text width=2.50cm,minimum height=1.15cm] (s0) at (0.0,0) {\textbf{1}\\Define 3 by 3 M};
\node[utpstep,text width=2.50cm,minimum height=1.15cm] (s1) at (2.95,0) {\textbf{2}\\Scan rows / columns};
\node[utpstep,text width=2.50cm,minimum height=1.15cm] (s2) at (5.9,0) {\textbf{3}\\Transpose / diagonal};
\node[utpstep,text width=2.50cm,minimum height=1.15cm] (s3) at (8.85,0) {\textbf{4}\\Check symmetry};
\draw[utparrow] (s0.east)--(s1.west);
\draw[utparrow] (s1.east)--(s2.west);
\draw[utparrow] (s2.east)--(s3.west);
\end{tikzpicture}
\end{adslabgraphic}
\end{adslabfigure}


\begin{utpsteps}\setcounter{enumi}{0}\item \textbf{Headers and initial definitions.}\end{utpsteps}

\texttt{matrix\_operations.cpp} --- lines 11--23.

\lstinputlisting[firstline=11,lastline=23]{all_experiments/Lab02/matrix_operations.cpp}


\tcbbreak

\begin{utpsteps}\setcounter{enumi}{1}\item \textbf{Implement rowSums, columnSums, transpose and helpers.}\end{utpsteps}

\texttt{matrix\_operations.cpp} --- lines 24--59.

\begin{adslabfigure}[Transpose: keep the value, exchange the indexes]
\begin{adslabgraphic}
\begin{tikzpicture}[font=\small]
\foreach \x/\a/\b/\c/\d/\e/\f/\g/\h/\i in {0/1/2/3/4/5/6/7/8/9,6/1/4/7/2/5/8/3/6/9}{
\foreach \val/\xx/\yy in {\a/0/0,\b/.85/0,\c/1.7/0,\d/0/-.68,\e/.85/-.68,\f/1.7/-.68,\g/0/-1.36,\h/.85/-1.36,\i/1.7/-1.36}{
\node[utpcell,minimum width=.8cm,minimum height=.62cm] at (\x+\xx,\yy) {\val};}}
\draw[utparrow] (2.4,-.68)--(5.3,-.68);
\node[utplabel,align=center] at (3.85,-.25) {\texttt{T[j][i] = M[i][j]}};
\node[utplabel] at (.85,.65) {M: original};
\node[utplabel] at (6.85,.65) {T: transposed};
\node[utplabel,align=center] at (3.85,-2.2) {M row sums: 6, 15, 24 \quad Diagonal sum: 15\\Maximum: 9 \quad M is not symmetric.};
\end{tikzpicture}
\end{adslabgraphic}
\end{adslabfigure}


\lstinputlisting[firstline=24,lastline=59]{all_experiments/Lab02/matrix_operations.cpp}


\tcbbreak

\begin{utpsteps}\setcounter{enumi}{2}\item \textbf{Connect the test cases in main.}\end{utpsteps}

\texttt{matrix\_operations.cpp} --- lines 60--99.

\lstinputlisting[firstline=60,lastline=99]{all_experiments/Lab02/matrix_operations.cpp}


\tcbbreak

\textbf{Build, run, and inspect matrix\_operations.}

From this program\textquotesingle s folder, run these commands in PowerShell:

\begin{lstlisting}[style=utpterminal,language={}]

g++ -std=c++17 -Wall -Wextra matrix_operations.cpp -o matrix_operations.exe

.\matrix_operations.exe

\end{lstlisting}

\textbf{Expected output excerpt.} Compare values, order, and pass/fail messages; spacing need not match.

\begin{lstlisting}[style=utpterminal,language={}]

M
    1    2    3
    4    5    6
    7    8    9

Row sums:
  row 0 sum = 6
  row 1 sum = 15
  row 2 sum = 24
Column sums:
  column 0 sum = 12
  column 1 sum = 15
  column 2 sum = 18

Transpose of M

\end{lstlisting}

The complete captured run is in \texttt{labs/outputs/matrix\_operations.txt}.

\begin{tryit}

Replace M with the identity matrix. Predict row/column sums of 1, diagonal sum of 3, maximum 1, and a true symmetry result.

\end{tryit}
